% ch3.tex -- Module 8 ESSENTIALS edition, Chapter 3: The quantum atom (Bohr, de Broglie,
% Schrödinger). Knowledge points only: no questions, no worked examples.
% check_numbers.py). The spectra are shown once, in Figure 1.1 (ch1:fig:spectra; the full
% book's ch3_spectra figure shows the same three spectra, so it is not repeated here). The
% diagrams are TikZ copied (some shrunk or rearranged) from the full chapter.
\chapter[The Quantum Atom]{The Quantum Atom: Bohr, de~Broglie and Schrödinger}\label{ch3:chap}
\covers{M8.IQ3.01, M8.IQ3.02, M8.IQ3.03, M8.IQ3.04, M8.IQ3.05}

\begin{summarybox}{In this chapter}
\begin{itemize}
  \item Why Rutherford's atom fails; Bohr's postulates and hydrogen's energy levels.
  \item The line spectrum: $E = hf$, $E = hc/\lambda$, the Rydberg equation, the Balmer series.
  \item De Broglie's matter waves, $\lambda = h/mv$, and the electron-diffraction evidence.
  \item The limits of Bohr's model, and what Schrödinger added.
\end{itemize}
\end{summarybox}

%% ==================================================================
\section{From Rutherford's atom to Bohr's energy levels}\label{ch3:sec:bohr}
\covers{M8.IQ3.01, M8.IQ3.03}

\begin{bigidea}
Classical physics says Rutherford's atom should collapse. Bohr fixed this with quantum rules:
the electron has only certain energies. It gives out or takes in light only when it jumps.
\end{bigidea}

\begin{keywords}
  \kw{Classical physics}{physics before quantum ideas: Newton's laws and Maxwell's
    electromagnetism.}
  \kw{Quantised}{only certain values are allowed.}
  \kw{Stationary state}{an allowed orbit, where the electron does not radiate (give out EM
    waves).}
\end{keywords}

\textbf{Why Rutherford's atom fails.} The \term{nuclear model} (1911) is a tiny,
dense, positive nucleus with electrons moving around it. The nucleus explained the gold-foil
results and is still accepted. But:
\begin{itemize}
  \item \textbf{It should collapse.} The orbiting electron accelerates $\to$ it gives out EM
        waves (Maxwell) $\to$ it loses energy $\to$ it spirals into the nucleus in about
        \qty{e-11}{s}. Real atoms are \term{stable}: they last for billions of years.
  \item \textbf{Wrong spectrum.} Spiralling in, it turns faster and faster, so its light would
        have every frequency: a \term{continuous spectrum}. Real hydrogen gives a
        \term{line spectrum}: a few sharp wavelengths. It also absorbs the same
        wavelengths (\cref{ch1:fig:spectra}).
  \item It does not say what the nucleus is made of, or how the electrons are arranged.
\end{itemize}

\textbf{Bohr's postulates (1913).} Bohr kept the nucleus and added \term{postulates}: rules
assumed without proof. Learn them by what they say (books number them differently):
\begin{enumerate}
  \item Electrons move in circular orbits, held by the electrostatic (Coulomb) attraction to
        the nucleus.
  \item Only certain orbits are allowed. In these \term{stationary states} the electron does
        \emph{not} radiate, so its energy stays constant.
  \item A jump between states emits or absorbs \emph{one} photon, with energy equal to the gap:
        $hf = E_\mathrm{i} - E_\mathrm{f}$ (emission; $E_\mathrm{f} - E_\mathrm{i}$ for
        absorption). i = initial, f = final. This is \textbf{conservation of energy}: the
        energy the atom loses is the energy of the photon.
  \item \term{Angular momentum} ($mvr$, the amount of turning motion) is quantised:
        $mvr = \dfrac{nh}{2\pi}$, $n = 1, 2, 3, \ldots$ (not on the formula sheet).
\end{enumerate}

\begin{figure}
  \centering
  \begin{tikzpicture}[x=0.9cm, y=0.8cm, font=\footnotesize]
    \foreach \y/\n/\E in {0/1/-13.6, 3.0/2/-3.40, 4.3/3/-1.51, 5.1/4/-0.85, 5.7/5/-0.54,
                          6.2/6/-0.38} {
      \draw[hscink, line width=0.8pt] (0,\y) -- (11.2,\y);
      \node[anchor=west] at (11.3,\y) {$n = \n$ \quad $\num{\E}$ eV};
    }
    \draw[hscink, line width=0.8pt, dashed] (0,6.8) -- (11.2,6.8);
    \node[anchor=west] at (11.3,6.8) {$n = \infty$ \quad 0 eV};
    \node[anchor=east, align=right] at (-0.1,0) {ground\\state};
    \node[anchor=east, align=right] at (-0.1,4.6) {excited\\states};
    \foreach \x/\y in {0.5/3.0, 0.9/4.3, 1.3/5.1, 1.7/5.7} {
      \draw[->, line width=0.9pt, hscfield!80] (\x,\y) -- (\x,0.04);
    }
    \node[draw=hscfield!80, rounded corners=1pt, align=center, fill=white, inner sep=2pt]
      at (3.7,0.65) {Lyman series\\(to $n=1$): ultraviolet};
    \foreach \x/\y/\col/\lab in {4.7/4.3/red!85!black/656, 5.5/5.1/cyan!75!blue/486,
                                 6.3/5.7/blue!80!black/434, 7.1/6.2/violet/410} {
      \draw[->, line width=1.1pt, \col] (\x,\y) -- (\x,3.04);
      \node[below, inner sep=1.5pt, text=\col] at (\x,3.0) {\lab};
    }
    \node[draw=hscink!60, rounded corners=1pt, align=center, fill=white, inner sep=2pt]
      at (6.1,1.95) {Balmer series (to $n=2$): 4 visible\\lines (in nm) + UV lines};
    \foreach \x/\y in {8.4/5.1, 8.8/5.7, 9.2/6.2} {
      \draw[->, line width=0.9pt, red!45!black] (\x,\y) -- (\x,4.34);
    }
    \node[draw=red!45!black, rounded corners=1pt, align=center, fill=white, inner sep=2pt]
      at (8.8,3.6) {Paschen series\\(to $n=3$): infrared};
    \draw[->, line width=0.9pt, hscink, dashed] (10.5,0.04) -- (10.5,6.75);
    \node[align=center, fill=white, inner sep=1pt] at (10.5,1.6)
      {ionisation\\from $n=1$:\\13.6 eV};
  \end{tikzpicture}
  \caption{Energy levels of hydrogen (spacing not to scale) and the three main series of
    lines.}
  \label{ch3:fig:levels}
\end{figure}

\begin{itemize}
  \item The allowed energies are \term{energy levels}, named by $n$, like the rungs of a ladder
        (\cref{ch3:fig:levels}). $n = 1$ is the \term{ground state}; higher levels are
        \term{excited states}. A jump between levels is a \term{transition}.
  \item Energies are negative: 0 eV is a just-free electron ($n = \infty$), and a bound one has
        less. Hydrogen: $E_n = -\dfrac{\qty{13.6}{eV}}{n^2}$ (not on the formula sheet).
  \item Only certain gaps exist, so only certain photon energies come out: a line spectrum.
        \textbf{Absorption} needs a photon that matches a gap \emph{exactly}.
  \item \textbf{Ionisation:} the electron reaches $n = \infty$ and leaves. The
        \term{ionisation energy} from $n = 1$ is \qty{13.6}{eV}; any extra becomes kinetic
        energy.
\end{itemize}

\begin{formulacard}{Photon energy}
  \cardformula{E = hf}
  \begin{symbols}
    \sym{$E$}{energy of one photon}{\si{J}}
    \sym{$h$}{Planck constant, \num{6.626e-34}}{\si{J.s}}
    \sym{$f$}{frequency of the light}{\si{Hz} ($\si{s^{-1}}$)}
  \end{symbols}
  \whentouse{a photon's energy from its frequency, or its frequency from an energy gap.}
  \trap{use the \emph{gap}, not one level's energy. Change eV to J first:
    $\qty{1}{eV} = \qty{1.602e-19}{J}$.}
\end{formulacard}

\begin{formulacard}[combine={$E = hf$ and $v = f\lambda$}]{Photon energy from wavelength}
  \cardformula{E = \frac{hc}{\lambda}}
  \begin{symbols}
    \sym{$E$; $\lambda$}{energy of one photon; wavelength of the light}{\si{J}; \si{m}}
    \sym{$h$; $c$}{Planck constant, \num{6.626e-34}; speed of light, \num{3.00e8}}{\si{J.s};
      \si{m.s^{-1}}}
  \end{symbols}
  \whentouse{wavelength $\leftrightarrow$ photon energy. Shorter $\lambda$: more energy.}
  \derive{for light $v = c$, so $f = \dfrac{c}{\lambda}$. Put it in $E = hf$.}
  \trap{nm to m: $\times 10^{-9}$. Not for particles with mass.}
\end{formulacard}


%% ==================================================================
\section{Hydrogen's line spectrum and the Rydberg equation}\label{ch3:sec:rydberg}
\covers{M8.IQ3.02, M8.IQ3.03}

\begin{bigidea}
Each line in hydrogen's spectrum is one kind of jump between two energy levels. The Rydberg
equation gives the wavelength of every line.
\end{bigidea}

\begin{keywords}
  \kw{Balmer series}{lines from jumps that end on $n = 2$: four visible lines, then UV.}
  \kw{Lyman / Paschen series}{jumps that end on $n = 1$ (UV, \qtyrange{91}{122}{nm}) /
    $n = 3$ (IR, \qtyrange{820}{1875}{nm}).}
\end{keywords}

\begin{itemize}
  \item Each element's lines are like a fingerprint (used to find hydrogen in stars). Balmer
        and Rydberg found formulas for the lines; Bohr explained them.
  \item Visible Balmer lines (\cref{ch1:fig:spectra}): \qty{656}{nm} (red, $3 \to 2$),
        \qty{486}{nm} (cyan, $4 \to 2$), \qty{434}{nm} (blue-violet, $5 \to 2$), \qty{410}{nm}
        (violet, $6 \to 2$). Then UV lines, down to the \term{series limit} (a jump from
        $n = \infty$), \qty{365}{nm}.
  \item \textbf{Link to the levels:} put $E_n = -\dfrac{\qty{13.6}{eV}}{n^2}$ into
        $\dfrac{hc}{\lambda} = E_\mathrm{i} - E_\mathrm{f}$. This gives the Rydberg equation,
        with $R = \dfrac{\qty{13.6}{eV}}{hc}$.
\end{itemize}


\begin{formulacard}{The Rydberg equation (hydrogen)}
  \cardformula{\frac{1}{\lambda} = R\left(\frac{1}{n_\mathrm{f}^2} - \frac{1}{n_\mathrm{i}^2}\right)}
  \begin{symbols}
    \sym{$\lambda$}{wavelength of the emitted (or absorbed) photon}{\si{m}}
    \sym{$R$}{Rydberg constant for hydrogen, \num{1.097e7}}{\si{m^{-1}}}
    \sym{$n_\mathrm{f}$; $n_\mathrm{i}$}{number of the final (end) level; of the initial
      (start) level}{no unit}
  \end{symbols}
  \whentouse{hydrogen lines only. Emission: $n_\mathrm{i} > n_\mathrm{f}$. Series limit:
    $n_\mathrm{i} = \infty$.}
  \trap{square each $n$ first. The equation gives $\frac{1}{\lambda}$: take the reciprocal at
    the end to get $\lambda$. $R$ is not resistance.}
\end{formulacard}

%% ==================================================================
\section{De Broglie's matter waves}\label{ch3:sec:debroglie}
\covers{M8.IQ3.04, M8.IQ3.01}

\begin{bigidea}
Light waves can act like particles. De Broglie (1924) proposed the reverse: particles can act
like waves. Electron diffraction confirmed it.
\end{bigidea}

\begin{itemize}
  \item \term{Wave--particle duality}: light and matter both show wave and particle behaviour.
  \item Photons carry \term{momentum} ($p$, amount of motion): $p = h/\lambda$ (not on the
        formula sheet).
  \item De Broglie's \term{hypothesis} (an idea to be tested): particles follow the same rule,
        with $p = mv$. So $\lambda = \dfrac{h}{mv}$.
  \item \term{Diffraction} (spreading out after a gap) is clear only when the gap is about the
        wavelength. Only waves diffract.
  \item A ball's $mv$ is large, so its $\lambda$ is far too small to notice. An electron at
        lab speeds has $\lambda \approx \qty{e-10}{m}$, about the atom spacing in a crystal. So
        crystals diffract electrons.
  \item \textbf{Evidence (1927):} Davisson and Germer, and G.\,P. Thomson
        (\cref{ch3:sec:experiments}). The evidence is that $\lambda$ from the pattern matched
        $\dfrac{h}{mv}$.
\end{itemize}

\begin{sidebyside}[0.38]
  \centering
  \begin{tikzpicture}[x=0.75cm, y=0.75cm]
    \foreach \n/\xx/\yy in {3/0/0, 4/3.0/0, 5/0/-2.9} {
      \begin{scope}[shift={(\xx,\yy)}]
        \draw[hscmuted!70, dashed, line width=0.5pt] (0,0) circle (0.75);
        \fill[hscforce] (0,0) circle (1.8pt);
        \draw[hscseriesA, line width=1.1pt] plot[domain=0:360, samples=240, smooth]
          ({(0.75+0.14*sin(\n*\x))*cos(\x)}, {(0.75+0.14*sin(\n*\x))*sin(\x)});
        \node[diagramlabel, align=center, anchor=north] at (0,-0.95) {$n = \n$};
      \end{scope}
    }
    \begin{scope}[shift={(3.0,-2.9)}]
      \draw[hscmuted!70, dashed, line width=0.5pt] (0,0) circle (0.75);
      \fill[hscforce] (0,0) circle (1.8pt);
      \draw[hscmistake, line width=0.8pt] plot[domain=0:720, samples=480, smooth]
        ({(0.75+0.14*sin(3.4*\x))*cos(\x)}, {(0.75+0.14*sin(3.4*\x))*sin(\x)});
      \node[diagramlabel, align=center, anchor=north] at (0,-0.95) {3.4 waves};
    \end{scope}
  \end{tikzpicture}
  \captionof{figure}{Electron waves around orbits (not to scale). Blue: allowed, a whole
    number of waves fits. Red: not allowed (two laps drawn, out of step).}
  \label{ch3:fig:orbitwaves}
\nextside
  \textbf{Standing waves explain Bohr's orbits} (\cref{ch3:fig:orbitwaves}).
  \begin{itemize}
    \item An allowed orbit fits a whole number of wavelengths: $2\pi r = n\lambda$ (not on the
          formula sheet). The wave joins itself: a \term{standing wave}.
    \item A standing wave does not travel, so it carries no energy away: ``no radiation''. In
          other orbits the wave cancels itself.
    \item Put $\lambda = h/(mv)$ into $2\pi r = n\lambda$: $mvr = nh/(2\pi)$.
          Bohr's angular-momentum postulate is \emph{explained}.
  \end{itemize}
\end{sidebyside}

\begin{formulacard}{de Broglie wavelength}
  \cardformula{\lambda = \frac{h}{mv}}
  \begin{symbols}
    \sym{$\lambda$}{de Broglie wavelength of the particle}{\si{m}}
    \sym{$h$}{Planck constant, \num{6.626e-34}}{\si{J.s}}
    \sym{$m$}{mass of the particle}{\si{kg}}
    \sym{$v$}{speed of the particle}{\si{m.s^{-1}}}
  \end{symbols}
  \whentouse{any moving particle with mass (electron, proton, neutron, ball). Not photons.
    Above about $0.1c$, use $\lambda = h/p_\mathrm{v}$ ($p_\mathrm{v}$ from the sheet).}
  \trap{faster means a \emph{shorter} wavelength and less diffraction. This $\lambda$ is not
    the decay constant.}
\end{formulacard}

%% ==================================================================
\section{Bohr's limits and Schrödinger's quantum model}\label{ch3:sec:schrodinger}
\covers{M8.IQ3.01, M8.IQ3.05}

\begin{bigidea}
Bohr's model works only for hydrogen-like atoms (one electron). Schrödinger (1926) treated the electron fully as a wave.
His model gives probabilities, not orbits: the basis of today's model.
\end{bigidea}

\textbf{Bohr's model.} Successes: stable atoms (by assumption) and hydrogen's line spectrum
(the Rydberg equation). Limitations:
\begin{enumerate}
  \item It works only for hydrogen and one-electron ions (such as He\textsuperscript{+}).
  \item It cannot explain why some lines are brighter than others (their intensity).
  \item It cannot explain \term{fine structure}: one line is really several very close lines.
        It cannot explain the \term{Zeeman effect} (lines split in a magnetic field) or the
        \term{Stark effect} (lines split in an electric field).
  \item Its rules are assumed, not explained. It mixes classical orbits with quantum rules.
  \item It gives the electron an exact path, which the uncertainty principle forbids.
\end{enumerate}

\Cref{ch3:tab:models} compares the three models.

\begin{table}
  \caption{Three models of the atom. The quantum numbers $l$ and $m_l$ give an orbital's shape
    and direction.}\label{ch3:tab:models}
  \centering
  \small
  \begin{tabular}{@{}L{21mm}L{36mm}L{38mm}L{42mm}@{}}
    \toprule
    & \colhead{Rutherford (1911)} & \colhead{Bohr (1913)} & \colhead{Schrödinger (1926)}\\
    \midrule
    Electron & a particle & a particle & a wave, described by $\psi$\\
    Where & any orbit & fixed circular orbits & orbitals: probabilities only\\
    Energy & any value & quantised by a postulate & quantised by the wave maths\\
    Stability & no: should collapse & yes, by assumption & yes: standing waves\\
    Quantum numbers & none & $n$ only (assumed) & $n$, $l$, $m_l$ (from the maths)\\
    Works for & $\alpha$ scattering only; no stable atom & hydrogen and one-electron ions
      only & all atoms
      (approximately)\\
    \bottomrule
  \end{tabular}
\end{table}

\textbf{What Schrödinger added:}
\begin{itemize}
  \item A \term{wave equation} for the electron in 3-D, built on de Broglie's matter waves. Its
        solutions are \term{wave functions}, $\psi$. Max Born (1926): $\psi^2$ gives the
        \term{probability} of finding the electron at each place.
  \item Heisenberg's \term{uncertainty principle} (1927): exact position and exact momentum
        cannot both be known. So there are no exact orbits.
  \item An \term{orbital} (1s, 2p, \ldots) is a 3-D cloud where \emph{one} electron is likely to
        be; it is not a path. In 1s the most likely distance is the \term{Bohr radius},
        \qty{5.29e-11}{m}.
  \item Energy levels and \term{quantum numbers} come out of the maths, not from postulates.
  \item It predicts line brightness and the simplest Zeeman splitting. Add electron spin
        (built-in turning) and Pauli's exclusion principle (no two electrons in an atom in the
        same state): it explains the periodic table. Add relativity (Dirac, 1928): fine structure.
\end{itemize}

%% ==================================================================
\section{Key experiments}\label{ch3:sec:experiments}
\covers{M8.IQ3.02, M8.IQ3.04}

\begin{experiment}{Investigating the Balmer series of hydrogen}
\textbf{Aim:} to measure hydrogen's visible lines and compare them with the Rydberg equation.
\textbf{What was done:} a hydrogen \term{discharge tube} (glowing gas) lit a grating
spectrometer (a diffraction grating, slit spacing $d$, with a scale to measure angles). Each
line's first-order angle $\theta$ gave $\lambda = d\sin\theta$. Then $\frac{1}{\lambda}$ was
plotted against $\left(\frac{1}{2^2} - \frac{1}{n_\mathrm{i}^2}\right)$. \textbf{Result:} four
lines (red, cyan, blue-violet, violet), each within 1\% of the Rydberg value; a straight line
through the origin, gradient \qty{1.10e7}{m^{-1}} $\approx R$. \textbf{Conclusion and
significance:} each line is a jump down to $n = 2$ (from $n = 3$ to 6): quantised energy
levels. These are the lines Bohr's model was built to explain. Astronomers use them to find
hydrogen in stars.
\end{experiment}

\begin{experiment}{Electron diffraction (1927): Davisson--Germer, Thomson}
\textbf{Aim:} to test de Broglie's matter waves. \textbf{What was done:} Davisson and Germer
fired \qty{54}{V} electrons at a nickel crystal in a vacuum; a detector counted them at each
angle. G.\,P. Thomson fired fast electrons through a thin metal film. \textbf{Result:}
Davisson--Germer: a peak at \ang{50} that moved with voltage; $\lambda = \qty{1.65e-10}{m}$ from
the pattern, $h/mv = \qty{1.67e-10}{m}$. Thomson: rings that shrank at higher voltage.
\textbf{Conclusion and significance:} electrons diffract, so they act as waves; $\lambda$
matched $h/mv$. This was the first direct evidence of matter waves. Electron microscopes use
the same idea.
\end{experiment}

\begin{summarybox}{Remember}
\begin{itemize}
  \item Rutherford's atom should collapse and give a continuous spectrum.
  \item Bohr: stationary states do \emph{not} radiate; a jump gives one photon,
        $hf = E_\mathrm{i} - E_\mathrm{f}$ (conservation of energy).
  \item Balmer series: jumps that end on $n = 2$; visible lines 656, 486, 434, \qty{410}{nm}.
  \item De Broglie: $\lambda = h/mv$, shown by electron diffraction; $2\pi r = n\lambda$.
  \item Schrödinger: $\psi^2$ = probability; orbitals, not orbits; levels from the maths.
\end{itemize}
\end{summarybox}
